\documentclass[11pt,reqno]{amsart}
\usepackage[T1]{fontenc}
\usepackage{lmodern,amssymb,booktabs}
\usepackage[margin=1.15in]{geometry}
\usepackage[colorlinks=true,linkcolor=blue,citecolor=blue,urlcolor=blue]{hyperref}
\newtheorem{theorem}{Theorem}
\newtheorem{proposition}[theorem]{Proposition}
\newtheorem{corollary}[theorem]{Corollary}
\newtheorem{lemma}{Lemma}[section]
\newcommand{\D}{\mathcal D}
\newcommand{\R}{\mathbb R}
\newcommand{\ord}{\operatorname{sort}_{\downarrow}}
\newcommand{\E}{\mathbb E}

\title[Seven-digit Kaprekar dynamics]{Seven-digit Kaprekar dynamics:\\ global law, two clocks, parity}
\author{Rustem Galyavin}
\thanks{Contact: \href{mailto:rustem@galyavin.com}{rustem@galyavin.com}.
ORCID: \href{https://orcid.org/0009-0006-6875-9170}{0009-0006-6875-9170}.}
\date{18 November 2026}
\subjclass[2020]{Primary 11A63; Secondary 37P99, 60F05}
\keywords{Kaprekar map, digit dynamics, arrival times, weak convergence}
\begin{document}
\begin{abstract}
Choose a nonconstant ordered word of seven base-$B$ digits uniformly, allowing
leading zeros. Under repeated Kaprekar subtraction let $T_B$ be its first arrival
time at a cycle. Along the bases of either parity, the normalized starting differences
and $T_B/B$ have a joint weak limit. At almost every starting point the
conditional time has two atoms. The faster time has probability $3/8$ in odd
bases and either $1/2$ or $3/4$ in even bases. Both marginal laws have support
$[0,1/4]$, and the even law has strictly smaller mean. We prove that every
arrival time is at most $54B+614132$ for $B\ge400$. All positive moments
and all fixed exponential moments converge along each parity. At first
arrival the normalized differences converge uniformly to $(3/4,1/2,1/4)$.
The proof combines an almost-everywhere entry theorem for a piecewise affine
map with exact integer returns and a computer-assisted potential inequality.
\end{abstract}
\maketitle
\section{Introduction}
In integer base $B\ge2$, arrange the seven digits of a word in decreasing
and increasing order and subtract the second number from the first.
Regard the result as another word of length seven. We write $\kappa_B$ for
this map and count time from the original word, so a word already on a cycle
has time zero. If $d_1\ge\cdots\ge d_7$ are its sorted digits, put
\[
 q=(u,v,w)=(d_1-d_7,d_2-d_6,d_3-d_5),\qquad X_B=q/B.
\]
Give each of the $B^7-B$ nonconstant ordered words equal probability. Let
$T_B$ be the first time $\kappa_B^{T_B}$ is on any cycle, and let $D_B$ be
the difference vector of that first cyclic word, divided by $B$.

\begin{theorem}[Global arrival law]\label{th:law}
Let $\D=\{(x,y,z):1>x>y>z>0\}$ and
$\rho(x,y,z)=5040(1-x)(x-y)(y-z)z$.
For $\beta=0,1$, as $B\to\infty$ with $B\equiv\beta\bmod 2$,
\[
 (X_B,T_B/B)\ \Longrightarrow\
 \rho(\xi)\,d\xi\,\bigl[p_\beta(\xi)\delta_{t_f(\xi)}
                    +(1-p_\beta(\xi))\delta_{t_s(\xi)}\bigr].
\]
The following construction defines the clocks and the probabilities for
Lebesgue-almost every $\xi\in\D$. Put
\[
 C(a,b,c)=\tfrac12(1+a,a+b,b+c),\qquad f(\xi)=C(\ord |2\xi_i-1|).
\]
In the ordered simplex $1\ge a\ge b\ge c\ge0$ put
\[
 K=\{2a\ge1,\ a+b\ge1,\ b+c\le a,\ b+2c\ge2a-1,\ 2b+c\le2-a\}.
\]
Iterate $f$ to the first visit to $C(K)$. Of that point and its next
two iterates exactly one, $(x,y,z)$, satisfies $x+z>2y$ and $x>y+z$.
Write $d=x-3z$. The clocks and weights at the original point $\xi$ are
\[
\begin{array}{c|cc|cc}
 &t_f&t_s&p_0&p_1\\ \hline
 d>0 &(9x-3z-6)/8&(5x-3z-3)/4&1/2&3/8\\
 d<0 &3(x+z-1)/4&(2x+6z-3)/4&3/4&3/8 .
\end{array}
\]
Almost everywhere $0<t_f<t_s<1/4$. Each marginal time law is absolutely
continuous with support $[0,1/4]$. The even-base limiting time is
stochastically smaller and has strictly smaller mean.
\end{theorem}

\begin{theorem}[Uniform capture and destination]\label{th:clock}
For every integer $B\ge400$ and every nonconstant seven-digit word,
\[
 T_B\le54B+614132,\qquad
 \|D_B-(3/4,1/2,1/4)\|_\infty\le 34/B .
\]
The same distance bound holds throughout each nonconstant cycle. The triple $(X_B,T_B/B,D_B)$ has the limit of Theorem~\ref{th:law}
with the constant destination appended. Write $\nu_\beta$ for the marginal
limiting time law. Along the same bases,
\[
 \begin{aligned}
 \E (T_B/B)^p&\longrightarrow\int t^p\,d\nu_\beta(t) &&(p>0),\\
 \E e^{\theta T_B/B}&\longrightarrow\int e^{\theta t}\,d\nu_\beta(t) &&(\theta\in\R).
 \end{aligned}
\]
In particular the normalized means and variances converge, and the
maximum arrival time among all words is $\Theta(B)$.
\end{theorem}

The conditional law in Theorem~\ref{th:law} is a statement about the joint
limit. At a given integer state the entire trajectory is deterministic.
Averaging over microscopic congruence classes near a starting point leaves
two possible times. The limit distinguishes the two base parities even after
the starting point is forgotten. For example, the point
$(0.84,0.51,0.24)$ is already on the positive canonical branch. There
the two times are $0.105$ and $0.12$. The shorter time receives half the
mass along even bases and three eighths along odd bases.

The proof has three parts. Section~3 shows that the real map $f$ almost
surely enters the region $C(K)$, where it is a rotation of order three
(Proposition~\ref{pr:entry}); the global law turns on this step.
Section~\ref{sec:facet} follows the exact integer map from there: three-step
returns carry the differences to the facet $2v=B$ and along it to the edge
$(x,1/2,1-x)$. Section~4 treats the final stretch near that edge, where the
time grows at one of two rates decided by residues modulo eight (Propositions~\ref{pr:edge}
and~\ref{pr:facet}). Section~6 transports these residues back to the start,
which gives the weights. Section~\ref{sec:uniform} proves the uniform bounds
with a potential that strictly decreases over nine steps while above $200$.

\subsection*{Earlier work}
Walden used symmetric differences in his node graph and constructed
seven-digit cycles near $(3B/4,B/2,B/4)$ for all large bases. He proved
that base four is the only base with a seven-digit globally attracting
fixed point and conjectured completeness of his list of cycles
\cite[\S3, Lemma 4.1, Theorem 4.2 and Conjecture 4.3]{Wa}.
Thakur gave another proof of the fixed-point result and observed the
shift $(3,2,1)$ in differences when the base increases by four
\cite[Theorem 4 and \S4.3]{Th}; on an explicit paired region of states
this shift is the seven-digit case of a covariance principle for every odd
number of digits \cite[Theorem~4]{G4}.
For every odd number of digits the continuous large-base map also decreases
the vector of successive spacings of the normalized differences, endpoint
spacings included, in majorization order \cite[Theorem~1]{G4}; the seven-digit
entry theorem below is proved separately. Hetmaniok, Pleszczy\'nski, Sobstyl and
Witu{\l}a tabulated the decimal seven-digit eight-cycle and its basin
\cite[entry $n=7$]{Het}; Nuez describes its $219$ difference classes
\cite[\S6]{Nu}. Kay and Downes-Ward studied fixed points and cycles for
any number of digits in odd and in even bases \cite{KD22,KD24}, and
Kahramanlar reduced the routine for any number of digits and any base to
the ordered opposite-position digit gaps \cite[\S2]{Kah} and computed the
determinants of its strict branches \cite[\S5]{Kah}, without seven-digit
asymptotics. Chen, Ono,
Schwartz and Thakur described four-digit dynamics in odd bases \cite{COST}.
For four digits in even bases, the difference
map yields cycle classifications \cite[\S2 and Theorem 1]{G1} and
arrival distributions \cite[\S2 and Theorems 1--2]{G2}. For five digits
a single normalized time at almost every starting point gives a common law
through all bases \cite[Theorem 1 and \S2]{G3}. Here the object is the
seven-digit joint distribution of position and first arrival as the base
grows. The proofs below require no classification of seven-digit cycles.

\section{Difference coordinates and input measure}
One subtraction of any word with differences $q$ produces the number
\[
 L_B(q)=(B^6-1)u+(B^5-B)v+(B^4-B^2)w .
\]
This is the seven-digit case of the gap-closure formula \cite[\S2]{Kah}.
The coefficients are positive, and each exceeds $(B-1)$ times the sum
of those to its right; thus $L_B$ is injective. Let $F_Bq$ be the differences
of $L_B(q)$ and $h_B(q)$ the first cyclic index for $F_B$. A cyclic word
has cyclic differences, and $L_B$ takes a cyclic difference vector to a
cyclic word. Therefore $h_B(q)\le T_B\le h_B(q)+1$.
More precisely, if the initial word is transient then $T_B=h_B(q)+1$:
for a transient difference vector its image under $L_B$ cannot be a
cyclic word, since a cyclic word has a cyclic predecessor and $L_B$ is
injective. This also shows that exactly one cyclic word has each cyclic
difference vector.
For $w>0$ column subtraction gives, from high digit to low,
\begin{equation}\label{eq:borrow}
 (u,v,w-1,B-1,B-w-1,B-v-1,B-u).
\end{equation}
Sorting this list and taking symmetric differences computes $F_B$,
including ties. At scale $B$ the corresponding real digits are $1$ and
the three pairs $(\xi_i,1-\xi_i)$. Sorting these pairs around $1/2$
gives the map $f$ in Theorem~\ref{th:law}. Let $\Gamma$ be a compact set
inside an open region on which the seven real digits have a fixed strict
order, so that $f$ is given there by a single sorting branch. On
$\Gamma$ the separations between the seven real digits are bounded below, so
for all sufficiently large $B$ they exceed the unit borrow corrections in
\eqref{eq:borrow} after scaling, and every integer $q$ with $q/B\in\Gamma$
follows the same branch. There $F_B$ is affine in $(q,B,1)$, differs from
$Bf(q/B)$ by a bounded vector, and has the linear part of $f$. This linear
part is integral of determinant $\pm1$: the diagonal factors $2$ in the
three signed radii cancel the factor $1/8$ in $\det C$. Compositions along
a fixed itinerary of such branches therefore permute the residues modulo
every integer bijectively (Section~6 carries out the induction). The
restriction to branches of the limiting map matters. A strict order of
the seven integer digits can have determinant $2$ where the limiting digits
tie: at $B=8k$ and $q=(6k,3k,2k)$ the seven digits are distinct, the branch is
$q\mapsto(B-w,2u-B,B-v-w-1)$ with determinant $2$, and the limiting
digits tie because $u+w=B$ (compare the odd central-axis restoration
theorem of \cite[\S5]{Kah}). Such walls are null.

Here is the population measure, including its lattice residues. For distinct
digits and prescribed $u>v>w>0$, there are $B-u$ choices of the smallest
digit. The two outermost adjacent gaps together sum to $u-v$, the next
two to $v-w$, and the inner two to $w$. A positive split of $g$ has $g-1$
choices. Each resulting sorted word has $7!$ orders. Words having two equal
digits number at most $\binom72 B^6$ in total. It follows by Riemann sums
that for every fixed Jordan measurable subset of $\D$, with a continuous
bounded weight, the limiting measure is $\rho(\xi)\,d\xi$. Within any one
residue class $q\bmod 8$ the limit is $1/512$ times this integral. The
same conclusion on the full simplex also shows $\int\rho=1$.

\section{Entry of the real map}
Define $H(a,b,c)=\ord(a,|a+b-1|,|b+c-1|)$. Direct substitution gives
$fC=CH$. The first coordinate of $H$ is $\max(a,1-b-c)$ and never
decreases. Each strict branch of $H$ has determinant of absolute value one:
before sorting its rows are $(1,0,0)$, $\pm(1,1,0)$ and $\pm(0,1,1)$.
Hence all finite inverse images of null sets are null. The initial map to
ordered radii has the same nonsingularity property. We may discard all
finite inverse images of the boundaries of branches, $K$, the planes
$a=1/2,1$, and the planes used for rotation and sign in the theorem.

\begin{proposition}\label{pr:entry}
Almost every point of the ordered radius simplex enters the interior of
$K$ after finitely many iterates of $H$.
\end{proposition}
\begin{proof}
Write $a_n$ for the first coordinate after $n$ steps. Below $a=1/2$,
one step either crosses the midpoint or gives $2a+b\ge1$. In this latter
region set $r=a-b$, $s=a+b+c-1$. By Lemma~\ref{l:entry}(i),
\[
 a'=a+(-s)_+,\qquad (r',s')=(|s|,r-s),\qquad t_+=\max(t,0).
\]
If $s$ stays nonnegative, the iterated matrix has rows $(0,1),(1,-1)$
and eigenvalues $\lambda=(\sqrt5-1)/2$ and $-1/\lambda$. Boundedness
on the simplex forces the expanding component to vanish, whence
$s=\lambda r$. Exclude this plane and all its finite inverse images.
Starting with $s=-t<0$, two further steps take $(r,t)$ to $(r+t,r)$
and increase $a$ by $t$, provided both steps remain below the midpoint.
This formula includes $t=0$. The Fibonacci recurrence from a nonzero
nonnegative pair makes the accumulated increments unbounded. Thus the
midpoint is crossed in finite time.

For $a\ge1/2$ measure the failure of the four remaining inequalities of $K$ by
\[
 W=(1-a-b)_++(b+c-a)_++(2a-1-b-2c)_++(2b+c+a-2)_+ .
\]
Write $Z$ for the union of the two segments with points $(a,2-3a,2a-1)$
and $(a,a,2a-1)$, $1/2\le a\le3/5$. Here $H$ interchanges the two indicated
points and $W=2a-1$. Exclude their finite inverse images also. We apply the
remaining parts of Lemma~\ref{l:entry}. If any $a_n>3/5$, stabilization
makes $a_n$ constant two steps later: apply the three-step equality at every
subsequent start. This constant is below one. On that plateau the six-step
decrement $\min(5a-3,1-a)$ is positive until $W$ reaches zero.

Otherwise $a_n\uparrow A\in(1/2,3/5]$ after passage above the midpoint.
Suppose the orbit avoids $K$. The sequence of six-step samples of $W$
decreases to a limit $w_0$. If $w_0=0$, the absorption inequality forces
entry since $\min(2a_n-1,1-a_n)$ is bounded away from zero. If $w_0>0$,
every accumulation point $p$ of the samples satisfies
$W(p)=W(H^6p)=w_0$. Equality places it among the two points of $Z$
over $A$. Compactness and continuity of the first five iterates give
$\operatorname{dist}(H^nq,Z\cap\{a=A\})\to0$ for the full orbit and
$W(H^nq)\to2A-1$. There are arbitrarily late visits near the first member
of the pair, taking an extra iterate from the second if necessary. The
neighborhood inequality in Lemma~\ref{l:entry} then supplies a later
state with $W<2a_n-1\le2A-1$, since crossing $3/5$ is excluded. Every third
subsequent value is at most this strictly smaller number. This contradicts
the full-orbit limit. Removing the finite inverse images of $\partial K$
turns entry into interior entry.
\end{proof}

On $K$ the sorted radius map is $H(a,b,c)=(a,1-b-c,a+b-1)$:
the ordering follows from $a+b+c\ge1$ and $a+2b+c\le2$. Substitution
in the inequalities gives $H(K)\subset K$ and $H^3=\mathrm{id}$ there.
On $C(K)$ the induced rotation is $R(x,y,z)=(x,x-z,y-z)$. It permutes
the three positive quantities $2y-1,\,2x-2y+2z-1,\,2x-2z-1$ cyclically
in the interior. The first is smallest at a unique rotation off their
equality planes. This is exactly the pair of strict inequalities in
Theorem~\ref{th:law}. The first $f$ step gives a point in the image of $C$,
so Proposition~\ref{pr:entry} applies to the original initial coordinates.
\section{Integer returns near the edge}
Consider the edge $(x,1/2,1-x)$ of the sorting regions. Measure distance
from it by the integer offsets $m=2v-B$, $S=u+w-B$.
\begin{proposition}[Edge clock]\label{pr:edge}
If $B\ge100000$, $|m|,|S|\le20$ and
$3B/4+1000\le u\le5B/6-1000$, there is a rate
$r\in\{3/2,2\}$ determined by $(B\bmod2,m,S)$ with
\[
 |h_B(q)-r(u-3B/4)|\le200 .
\]
\end{proposition}
\begin{proof}
A three-step return sends the two offsets to new offsets and shifts $u$
by a constant. To compute the transition for any reachable offsets use
\eqref{eq:borrow} three times at $B=100000+(m\bmod2)$, $u=80000$.
Lemma~\ref{l:edge} makes this calculation valid over the entire strip
$3B/4+128\le u\le5B/6-128$. The initial path reaches one of the following
offset cycles; each entry shows offsets $(m,S)$ in cyclic order. Write
$\ell$ for the number of three-step returns, and $D$ for their total
decrease of $u$.
\[
\begin{array}{c|l|rr}
 B\bmod2&\text{cycle}&\ell&D\\\hline
0&(0,-2),(4,2)&2&3\\
0&(2,-1),(6,1),(2,0)&3&6\\
0&(4,0)&1&2\\
1&(-1,-3),(3,3),(-1,-1),(3,1)&4&6\\
1&(-1,-2),(3,2)&2&3\\
1&(1,-2),(5,2)&2&3\\
1&(1,-1),(5,1)&2&4
\end{array}
\]
Set $r=3\ell/D$. The transient endpoints and one further circuit move
$u$ by at most $30$ from its initial value, with intermediate excursions
at most $20$ in a block. After the transient, repeat whole circuits until
the first endpoint below $H_B=\lceil3B/4\rceil+256$. It lies in
$[H_B-D,H_B-1]$. At any point within a circuit the displacement from its
start has absolute value at most $2\cdot30+20=80$. Since circuit starts
decrease, the margins at initial entry and at the threshold justify every
return used. Every skipped word in difference coordinates has $u-3B/4\ge128$.

For every possible endpoint, direct subtraction followed by affine lifting
gives the rest of the path and its first cyclic index $h_{\rm end}$
(Lemma~\ref{l:edge}). Its cycle has $u-3B/4<128$ throughout. Hence the
first capture belongs to the final path. If the initial transient takes
$k$ blocks and shifts $u$ by $e$, the time minus $r(u-3B/4)$ is
\[
 3k+re+h_{\rm end}-r(u_{\rm end}-3B/4).
\]
The interval in Lemma~\ref{l:edge} bounds all these discrepancies.
\end{proof}

\section{From the period-three region to the edge}\label{sec:facet}
With exact digits, the interleaving order of the seven entries in
\eqref{eq:borrow} (indices $0,\ldots,6$) is $(3,0,4,1,5,2,6)$ on
one sorting region. Its step is $(u-1,u-w+1,v-w)$.
Three successive such steps translate $q$ by $\Delta=(-3,-2,-1)$.
In particular the three inner gaps $(u-v,v-w,w)$ themselves change
in one step to $(w-2,u-v+1,v-w)$ there. If an itinerary starts with
these three orders, the exact number of complete repetitions is
\[
 \min_{g_{\rm lin}(\Delta)<0}
 \left(1+\left\lfloor g(q)/[-g_{\rm lin}(\Delta)]\right\rfloor\right).
\]
Here $g$ ranges over all input and pulled-back sorting inequalities
$g\ge0$ of that three-step word, including $w\ge1$; at a tie the
smaller index comes first, so a strict comparison has integer margin one.
The formula follows by evaluating $g(q+j\Delta)$ for consecutive $j$.
The bound on $w$ supplies a negative slope. For the global law we also
need the change of sorting order at the facet $2v=B$.

Besides $S,m$, use the forms
\[
 \begin{gathered}
 E=2u-2w-B=S+d,\qquad d=u-3w,\\
 J=B-u-v+2w,\qquad K_1=B-2u+2v-w,\qquad L=B-v-w .
 \end{gathered}
\]
Consider integer states satisfying
\begin{align}
 &B\ge100000,\quad u\le B-1000,\quad w\ge1000,\quad m\ge0,\notag\\
 &S-m,\ E-m,\ J,\ K_1,\ L\ \ge1000. \label{eq:source}
\end{align}
The positive source also requires $d\ge8000$, $3d\le B-24000$; the
negative source requires $-d\ge8000$, $B+6d\ge24000$.

\begin{proposition}[Facet transit]\label{pr:facet}
For any original word with differences in either source,
$|T_B-Bt_r(q/B)|\le216$, where $t_r$ is the clock of the corresponding
sign in Theorem~\ref{th:law}, evaluated directly there: $r=3/2$ selects
$t_f$ and $r=2$ selects $t_s$. The rate is determined by the residue
table below.
\end{proposition}
\begin{proof}
Lemma~\ref{l:facet} gives, under its weaker guards, the three-step
update
\[
 (u,w,m)\mapsto(u-3,w-1,|m+1|-5).
\]
Let $m_j$ denote the last coordinate after $j$ such returns. Starting at
$m\ge0$ it decreases by four until entering $-4\leftrightarrow-2$
(even base), or $-3$ or $-5\leftrightarrow-1$ (odd base), in at most
$\lfloor m/4\rfloor+1$ steps. Thus $m_j\ge m-4j$, $-5\le m_j\le m$.
Both $S$ and $E$ fall by four each time and $d$ is fixed. For $d>0$
take $N=\lfloor(S-16)/4\rfloor$, $s=S-4N\in[16,19]$. For $d<0$ take
$N=\lfloor(E-12)/4\rfloor$, $e=E-4N\in[12,15]$. The margins in
\eqref{eq:source} put $m_N$ on its cycle. At every $j<N$ both $S_j,E_j$
are at least $16$. Before the bounded phase $S_j-m_j,E_j-m_j$ are
constant, and in that phase $m_j\le0$. Furthermore
\[
 J_j-J=L_j-L=j+(m-m_j)/2,\qquad (K_1)_j-K_1=7j+m_j-m\ge3j .
\]
The first coordinate never exceeds its starting value. The final $w_N$
in the formulas below exceeds $B/6$; all earlier values are larger.
These facts verify the guards throughout, so the elapsed time is exactly
$3N$ up to this point.

For $d>0$ the endpoint has
\[
 u_N=(3B+d+3s)/4,\quad w_N=(B-d+s)/4,\quad (m,S)=(m_N,s).
\]
It satisfies Proposition~\ref{pr:edge} by the two bounds on $d$.
For $d<0$ the endpoint instead has
\[
 u_N=(3B+3e-2d)/4,\qquad w_N=(B+e-2d)/4 .
\]
One more subtraction has first coordinate $u_N-1$ and offsets
\begin{equation}\label{eq:neg}
 m'=e+2,\qquad S'=
 \begin{cases}(e+m_N)/2-1&m_N\ge-1,\\
               (e-m_N)/2-2&m_N\le-2.
 \end{cases}
\end{equation}
Lemma~\ref{l:facet} places this image in the edge strip. In both cases
the edge clock is positive, so no preceding difference state was cyclic.
Adding the elapsed time and the edge clock, with the possible extra step
from $h_B$ to $T_B$, gives $T_B$ as the leading term
\[
 \tfrac34S+\tfrac r4d\quad(d>0),\qquad
 \tfrac34E-\tfrac r2d\quad(d<0),
\]
which equals $Bt_r(q/B)$ in both cases, plus the edge error and the
bounded terms $3s(r-1)/4+\epsilon$ and $3e(r-1)/4+1-r+\epsilon$,
respectively, where $0\le\epsilon\le1$. These terms have absolute value
below $16$, and the edge error is at most $200$.

For completeness, the following table lists the fast columns. Columns
are indexed $0,1,2,3$ for the four increasing values of $m\bmod8$
compatible with the base parity. A dash means the empty set. The row
variable is $S\bmod8$ on the positive side and $E\bmod8$ on the
negative side. A dot indicates a parity-incompatible row.
\[
\begin{array}{c|cc|cc}
 &\multicolumn{2}{c|}{d>0}&\multicolumn{2}{c}{d<0}\\
 \text{row}& B\text{ even}&B\text{ odd}&B\text{ even}&B\text{ odd}\\ \hline
 0&12&02&0123&\cdot\\
 1&03&02&\cdot&-\\
 2&03&-&03&\cdot\\
 3&03&02&\cdot&023\\
 4&03&02&0123&\cdot\\
 5&12&02&\cdot&-\\
 6&12&-&12&\cdot\\
 7&12&02&\cdot&012
\end{array}
\]
Indeed, the residue of $S$ or $E$ determines its endpoint remainder
and the parity of $N$. The eventual value $m_N$ depends only on that
parity and $m\bmod8$: increasing $m$ by eight adds two initial
subtractions of four. Evaluating the indicated edge graph at $(m_N,s)$
or \eqref{eq:neg} gives the table (Lemma~\ref{l:facet}).
\end{proof}

\subsection*{Which sources are reached}
In coordinates $\xi=(x,y,z)$ put $a=2x-1$, $b=2y-a$,
$c=2z-b$. Substitute these in the ordered simplex, the inequalities of
$K$, and the two inequalities selecting the canonical rotation.
The resulting closed polyhedron is described by the homogeneous limits
of \eqref{eq:source} (with $B=1$ and all additive margins removed),
ordered normalized coordinates, and the relevant sign of $x-3z$.
Its vertices on either side are given by the two lists
\[
 \begin{array}{c|l}
 + & A,\ U,\ G,\ D,\ V\\
 - & A,\ G,\ F,\ V
 \end{array}
 \quad
 \begin{aligned}
 A&=(3/4,1/2,1/4), & U&=(5/6,1/2,1/6),\\
 G&=(6/7,1/2,2/7), & D&=(9/10,1/2,1/5),\\
 F&=(5/6,1/2,1/3), & V&=(1,2/3,1/3).
 \end{aligned}
\]
The vertex calculation is specified in Lemma~\ref{l:facet}. In particular
$0<d/B<1/3$ on the interior of the positive side and
$-1/6<d/B<0$ on the negative side. Thus every compact set in a
canonical source interior satisfies all integer buffers after scaling by
sufficiently large $B$, even with bounded additive errors in its coordinates.

\section{Passage to the population law}
Partition the full-measure set of strict entry points by the finite
sequence of sorting branches needed to enter $C(K)$, by the subsequent
rotation, and by sign. Refine along the finitely many relevant comparison
planes in each fixed itinerary. This gives countably many open polyhedra
up to null boundaries. On any compact polytope inside one, the rescaled
digits along the itinerary have strict separations bounded below. By
induction over the fixed itinerary, as in Section~2, the exact map for all
sufficiently large $B$ follows that prefix. Its
endpoint is $B$ times the canonical image plus a bounded offset depending
only on the branch sequence. It lies in \eqref{eq:source} with the
prescribed sign. The positive remaining time there justifies adding the
prefix length. For example the positive remaining time follows already
from the transient edge reached in Proposition~\ref{pr:facet}. Both the
prefix length and the discrepancy in the unnormalized affine clocks are
bounded on this compact set.

Fix $B\bmod8$. The affine prefix permutes all $512$ input residues, and
each has the same limiting conditional frequency in any subpolytope.
On the positive side, $(u,v,w)\mapsto(S,m)\bmod8$ has $16$ preimages for
each of its $32$ admissible outputs. On the negative side,
$(u,v,w)\mapsto(E,m)\bmod8$ has $32$ preimages for each of $16$ outputs.
Indeed $u+w$ has eight representations modulo eight, while $2v$ has
two; for $2(u-w)$ there are sixteen representations per admissible
value. The table consequently gives $16/32$ and $12/32$ on the positive
side, $12/16$ and $6/16$ on the negative side. These proportions depend
only on $B\bmod2$.

For each bounded continuous joint test, the residue Riemann sums give the
asserted integral on the compact polytope: the affine clocks are continuous
there, and the normalized error vanishes uniformly. Exhaust the entry set
by finite unions of such interior polytopes, using disjoint interiors.
The limiting and discrete omitted masses can both be made arbitrarily small
by the input Riemann sum. The size of the test bounds their contribution,
regardless of entry lengths outside the chosen polytopes. This proves the
weak convergence along each residue modulo eight and hence, by the equality
of the four limits, along each parity.

The two times differ by $d/(8B)$ or $-d/(4B)$ in the notation of the
integer facet (equivalently by $(x-3z)/8$ or $(3z-x)/4$ in canonical
coordinates). Both affine forms are nonconstant on every entry cell
since the prefix is invertible there. Hyperplane slicing gives an absolutely
continuous pushforward of the restricted density on each cell, and hence
of their countable sum. Evaluating at the listed vertices gives the
range $[0,1/4]$ as an upper and lower bound, strict in the interiors.
On the segment $AV$ both clocks increase linearly from $0$ to $1/4$.
Every open subinterval is reached on an open set by perturbation into
the positive interior; this set has positive $\rho$-mass. This proves
the support assertion. Coupling the two laws using the same initial
point and a common uniform choice of atom orders them, since $p_0>p_1$.
Their mean difference (odd minus even) is
$\int_{\D}\rho(p_0-p_1)(t_s-t_f)>0$.
This completes the proof of Theorem~\ref{th:law}.
\begin{corollary}[The period-three region and its preimage]\label{cor:regional}
Let $Q=C(K)$ and $P=f^{-1}(Q)$ in the ordered initial simplex. For
$S=Q$ or $P$, conditioning on $X_B\in S$ gives the law of
Theorem~\ref{th:law} restricted to $S$ and divided by $M_S=\int_S\rho$,
where
\[
 M_Q=1339759/43200000,\qquad M_P=937297/5400000 .
\]
\end{corollary}
\begin{proof}
Both boundaries have zero limiting mass, so restriction follows by weak
convergence. To compute $M_Q$, take the two positive tetrahedra
$(V,A,U,D),(V,A,D,G)$ and the negative tetrahedron $(V,A,G,F)$
from Section~\ref{sec:facet}. Their three rotations partition $Q$ up to faces.
The eight inverse branches of $f$ send $C(a,b,c)$ to the decreasing sort
of $( (1\pm a)/2,(1\pm b)/2,(1\pm c)/2 )$, with signs chosen
independently. Each is affine on the ordered cone and their interiors are
disjoint. Their images of the nine tetrahedra cover $P$. The two
integrals follow by the rational calculation in Lemma~\ref{l:facet}.
\end{proof}

\section{Uniform bounds on the integer trajectory}\label{sec:uniform}
A positive state in this section means $w>0$. For $B\ge18$ a positive
start either stays positive or reaches zero in the next two steps. Indeed,
three consecutive positive states cannot have a zero successor by
Lemma~\ref{l:zero}(i); apply this to the last three states before a
putative later first zero.

For a zero start $(u,v,0)$ with $v>0$ write $B=2b+\beta$.
Lemma~\ref{l:zero}(ii) gives a positive image. A second-step return
to zero occurs only from $(b,b,0)$ in even bases and from
$(b,b,0),(b+1,b,0),(b+1,b+1,0)$ in odd bases. Its image at that
second step is $(b+1,1,0)$ in all cases. Otherwise the second and third
images are positive. The first-access bound applied to the first image
now excludes any subsequent zero in this latter case. The exceptional
return point lies outside the list for $b\ge9$, so it leaves permanently
on the next step. Finally, the exact first image of $(u,0,0)$ is
\[
 (\max(B-u,u-1),\min(B-u,u-1),0).
\]
When $u=1$ two steps give $(B-2,1,0)$, which leaves permanently on the
next step. In the other cases the middle difference of the image is
positive and its first two coordinates sum to $B-1$. Among the exceptional
sources this can only be $(b,b,0)$ in an odd base. Thus any zero start
is permanently positive by step four; an arbitrary start is permanently
positive by step six.

For $g=(B-u,u-v,v-w,w)$ introduce
\[
 V(q)=\sum_{i=1}^4|4g_i-B|,\qquad X=(4u-3B,4v-2B,4w-B).
\]
The gaps are nonnegative with sum $B$, giving $0\le V\le6B$ by convexity
on the simplex. The four centered gaps are $-X_1,X_1-X_2,X_2-X_3,X_3$.
Since these sum to zero, every partial sum has absolute value at most
half their total absolute sum. Thus $\|X\|_\infty\le V/2$.

On any permanently positive orbit, Lemma~\ref{l:clock} gives
$V(F_Bq)\le V(q)+8$ and a strict decrease at time nine whenever $V(q)>200$.
At times $0,9,18,\ldots$ integer descent yields a sample with $V\le200$
within $54B$ steps. All subsequent sample values are at most $272$:
a sample at most $200$ gains at most $72$ and a larger one strictly
decreases. Intermediate values are therefore at most $336$. From this
point on $X$ lies in $[-168,168]^3$, with each coordinate in one fixed
residue class modulo four. There are at most $85^3=614125$ possible
states there. Determinism forces a repeat within that many steps.
Adding six initial steps and at most one to pass to the original word
gives the bound of Theorem~\ref{th:clock}.

On a cycle take a maximum of $V$ and go nine steps backwards on the
same cycle. If the predecessor has $V>200$, its value strictly exceeds
the maximum, an impossibility. Otherwise the maximum is at most $272$.
Thus $\|X\|_\infty\le136$ on every cycle. The differences of the first
cyclic word are cyclic, yielding the asserted uniform destination bound.
This also shows that the destination can be appended to the weak limit.
The family $T_B/B$ is uniformly bounded for $B\ge400$, so every fixed
power $p>0$ and every exponential with fixed real parameter is a bounded
continuous test on the common range. Finally, the support of each limiting
law gives $\liminf_{B\to\infty}\max T_B/B\ge1/4$ by testing the open set
$(1/4-\epsilon,\infty)$ along both parities and then decreasing $\epsilon$.
This proves Theorem~\ref{th:clock}.

\section{Computations and verification}
The accompanying directory runs from a fresh copy with
\texttt{bash reproduce.sh}. It regenerates the real and integer solver
queries, checks rational combinations of sorting inequalities exactly,
and follows terminal paths by digit subtraction. The appendix specifies
the finite assertions used in the proof. Z3 5.4.1 and cvc5 1.4.2 were
used for real entry and zero-face assertions. Z3 5.4.1 and Yices 2.7.0
were used for the potential. The two unbounded longitudinal-strip
checks used Z3 5.4.1. Python 3.9.6 with exact fractions performs the
rational and finite graph replays. Numerical checks additionally use
SymPy 1.14.0 and a C++17 compiler.

The independent column-subtraction implementation visits all $22\,315\,002$
nonconstant difference vectors in bases $129,130,400,401$. The maximal
observed difference arrival times are $56,49,113,127$, respectively.
Whole-integer subtraction also checks twelve trajectories selected by
their finite leading prefixes, totaling $1\,643\,929$ digit updates.
Exact weighted distributions in bases $96,97,192,193$ account for all
$B^7-B$ words at each base. Finite examples check the implementations.
The conclusions for unbounded bases use the appendix and the preceding
analytic arguments. The time integrals in the theorem specify limiting
moments without a numerical quadrature claim. Neither a sharp maximum nor
an exact list of all cycles is needed here.

An independent re-derivation of every reduction from the text of an earlier
draft of this paper was also carried out, with separately written encodings.
Z3 found no counterexample to the entry, facet, zero-face and potential
assertions, with no upper bound on $B$, and Yices confirmed the entry
assertions. The edge graphs, the $140$ terminal families and both regional
masses were recomputed independently and agree exactly. Word-level
subtraction at the bases $200000$--$200007$, $1000003$ and $1000004$ placed
all $440\,320$ sampled arrival times within $99$ steps of one of the two
predicted clocks. In each box of $512$ words, one per residue, the fast
share was $1/2$ on the positive side and $3/4$ on the negative side in even
bases, and $3/8$ on both sides in odd bases. A separate exhaustive
enumeration at the bases $18$--$300$, $400$, $401$, $500$--$507$, $1000$
and $1001$ found every nonconstant cyclic state within $3.5/B$ of
$(3/4,1/2,1/4)$ in the sup norm. These encodings and their logs are part of the ancillary
material.

\appendix
\section{Finite assertions and their encodings}\label{sec:appendix}
In the following lemmas an integer state always satisfies
$B>u\ge v\ge w\ge0$, $u>0$. On reals write $q_j=H^j(a,b,c)$ and
$a_j=(q_j)_1$. Real absolute values, maxima, minima and sorting use
piecewise affine expressions with non-strict comparisons at equality.

\begin{lemma}[Eight real inequalities]\label{l:entry}
For $0\le c\le b\le a<1$ the following assertions hold.
\begin{enumerate}
\item[(i)] If $a,a_1<1/2$, then $2a_1+(q_1)_2\ge1$. If $a<1/2$
and $2a+b\ge1$, then, with $s=a+b+c-1$,
$q_1=(a+(-s)_+,a-s_+,1-a-b)$.
\item[(ii)] For $a\ge1/2$, $W(q_3)\le W(q_0)$.
\item[(iii)] For $a>3/5$, $a_3=a_2$. If additionally $a_6=a$,
then $W(q_6)\le\max(0,W(q_0)-\min(5a-3,1-a))$.
\item[(iv)] If $a>1/2$, $a_6<1$ and $W(q_6)=W(q_0)>0$, then $q_0\in Z$.
\item[(v)] If $a>1/2$ and $W(q_0)<\tfrac12\min(2a-1,1-a)$,
then $W(q_6)=0$.
\item[(vi)] If $1/2<a\le3/5$, $(b,c)\ne(2-3a,2a-1)$ and
$\max(|b-2+3a|,|c-2a+1|)<(2a-1)/100$, then some $0\le j\le4$
has $a_j>3/5$ or $W(q_j)<2a_j-1$.
\end{enumerate}
\end{lemma}
\noindent
\texttt{seven\_entry\_audit.py} takes a triple of real variables, forms
six iterates, and negates the eight claims in quantifier-free linear real
arithmetic. A sorted triple is encoded by its maximum, its sum less the
maximum and minimum, and its minimum. For the stabilization query it
suffices to exclude $a_3>a_2$, by monotonicity. Each negation is UNSAT in
both Z3 and cvc5.

\begin{lemma}[Edge graph and terminal lifts]\label{l:edge}
For $\beta=0,1$ start with all offsets $|m|,|S|\le20$,
$m\equiv\beta\pmod2$, and close under the representative three-step rule
in Proposition~\ref{pr:edge}. On every reachable node, for all bases
$B\equiv\beta\bmod2$, $B\ge100000$, and every integer $u$ with
$3B/4+128\le u\le5B/6-128$, the representative rule gives the exact
three-step offsets and displacement in $u$. Every partial step of a block
moves $u$ by at most $20$ from its value at the start of the block.
The graphs have $884,845$ nodes and precisely the cycles listed in that
proposition. Starting in the indicated offset square, the transient lengths
in blocks are at most $10,11$. The cumulative displacement at each
block endpoint through the transient and one circuit has absolute value
at most $24,30$, respectively.

For each reachable cycle phase, compatible $B\bmod4$, and
$u_{\rm end}=H_B-j$ with $1\le j\le D$, the whole path from the endpoint
to first repetition is the translate of its path in the representative
base $100000+(B\bmod4)$ by $(3,2,1)(B-100000-(B\bmod4))/4$.
Every terminal cycle on these paths has $u-3B/4<128$. The discrepancy
$3k+re+h_{\rm end}-r(u_{\rm end}-3B/4)$ over all admissible initial
offsets and terminal choices lies in $[-43/2,343/8]$.
\end{lemma}
\noindent
In \texttt{seven\_boundary\_audit.py} write $B=2b+\beta$ and start at
$(u,b+(\beta+m)/2,2b+\beta-u+S)$. Sort the seven affine forms in
$(u,b,1)$ at the representative; retain every comparison and each
intermediate state-validity inequality. A reversed index in a tie-broken
comparison requires a margin of one. Composing the three steps gives
constants for both offsets and all partial first-coordinate displacements.
The disjunction that any of $328$ or $330$ distinct guards fails on the
given strip is UNSAT in Z3 for each parity. Traversing the resulting finite
graphs from $861$ even and $820$ odd initial nodes yields the cycle and
excursion data. For the $140$ endpoint families, compute up to first
repetition at the representative and check descending order of both digit
values and their slopes $(3,2,1,4,3,2,1)$ under $B\mapsto B+4$. The
borrowed digit comparisons thus persist for all nonnegative shifts.
The symmetric slope differences give $(3,2,1)$ at every step. Translation
preserves all equalities and inequalities between entire states, and hence
the index of the first repetition. The finite enumeration of the discrepancy
uses these first cyclic indices.

\begin{lemma}[Facet returns, exits and residue tables]\label{l:facet}
Use $m,S,E,d,J,K_1,L$ as in Section~\ref{sec:facet}. If
\[
 \begin{gathered}
 B\ge100000,\quad u\le B-1000,\quad w\ge1000,\quad m\ge-6,\\
 S,\ S-m,\ E,\ E-m\ge16,\qquad J,\ K_1,\ L\ge8 ,
 \end{gathered}
\]
then three steps give $u'=u-3,\ w'=w-1,\ m'=|m+1|-5$.
For $B\ge100000$, $-d\ge8000$, $B+6d\ge24000$, $E=e\in\{12,13,14,15\}$
and $m\in\{-5,-4,-3,-2,-1\}$ of the same parity as $e$, the next state
has $u'=u-1$, offsets \eqref{eq:neg} with $m_N=m$, and meets the full
domain of Proposition~\ref{pr:edge}.
The closed canonical sign regions have precisely the vertices listed
in Section~\ref{sec:facet}, and applying the edge rule to each endpoint residue gives
precisely the fast columns in that section. Integrating $\rho$ over the
nine tetrahedra for $Q$ and their $72$ inverse images gives respectively
the two rational numbers in Corollary~\ref{cor:regional}.
\end{lemma}
\noindent
\texttt{seven\_return\_checks.py} reconstructs thirteen three-step
sorting words ($m=-6,\ldots,5$ and $m\ge6$) and ten one-step exit words.
Each retained affine guard is supplied as a nonnegative rational linear
combination of the hypotheses, branch constraints and $1\ge0$. The
replay reconstructs all rows and final matrices, checks $490$ identities
over $\mathbb Q$, and rejects missing rows, changed coefficients or a
changed domain. Vertex checks intersect every independent triple of
bounding planes and retain exactly the feasible points. The replay does
this for the limiting integer source domains (including their sign
bounds); \texttt{seven\_tables.py} repeats it starting directly from the
defining inequalities of the canonical $C(K)$ regions. For the positive
table one evaluates the bounded phase $m_N$ at residues of $(S,m)$; for
the negative table one evaluates \eqref{eq:neg} at residues of $(E,m)$.
The edge graph supplies the rate at each resulting offset. Both the column
lists and their counts are checked. For the integrals, expand the four
linear factors of $\rho/5040$ in barycentric coordinates. For a tetrahedron
with affine Jacobian $J_T$, each degree-four monomial with exponents
$k_0,\ldots,k_3$ has integral $|J_T|\prod k_i!/7!$. Summing these exact
terms gives the masses. The older, restricted positive source
also has a separate $338$-identity replay in \texttt{seven\_facet\_checks.py};
the two-sided return above suffices for the theorem.

\begin{lemma}[First access and zero-origin returns]\label{l:zero}
For every integer base $B\ge18$ the following hold.
\begin{enumerate}
\item[(i)] If $q,F_Bq,F_B^2q$ are positive then $F_B^3q$ is positive.
\item[(ii)] For $q=(u,v,0)$, $v>0$, $F_Bq$ is positive. The second
image lies on $w=0$ exactly at the sources listed in Section~\ref{sec:uniform}, and then
equals $(\lfloor B/2\rfloor+1,1,0)$. If the second image is positive,
so is the third.
\item[(iii)] For $q=(u,0,0)$ the first image is the sorted pair
$\max(B-u,u-1),\min(B-u,u-1)$ followed by zero.
\end{enumerate}
\end{lemma}
\noindent
For (i) \texttt{seven\_first\_access.py} encodes three positive sources
by \eqref{eq:borrow} and asks for a final $w=0$ with $B\ge18$.
For (ii) the borrowed digits are
$(u,v-1,B-1,B-1,B-1,B-v-1,B-u)$, and for (iii) they are
$(u-1,B-1,B-1,B-1,B-1,B-1,B-u)$. The routine
\texttt{seven\_zero\_transience\_audit.py} writes $B=2b+\beta$,
$b\ge9$, and negates the three assertions in (ii) separately in each
parity. It checks the formula in (iii) for every $B\ge18$. These seven
and the first-access negation are UNSAT in both Z3 and cvc5. Sorting uses
the same network as the next lemma.

\begin{lemma}[Potential inequalities]\label{l:clock}
For every integer $B\ge2$ and every positive source,
$V(F_Bq)\le V(q)+8$. If all ten states $q,F_Bq,\ldots,F_B^9q$ are
positive and $V(q)>200$, then $V(F_B^9q)<V(q)$.
\end{lemma}
\noindent
\texttt{seven\_clock\_audit.py} removes one maximal digit $B-1$ from
\eqref{eq:borrow} and sorts the remaining six using the following
comparators in order, placing the larger digit in the lower index:
\[
 (1,2),(4,5),(0,2),(3,5),(0,1),(3,4),
 (2,5),(0,3),(1,4),(2,4),(1,3),(2,3).
\]
For the resulting list $d_0,\ldots,d_5$, the image is
$(B-1-d_5,d_0-d_4,d_1-d_3)$. The network is checked on every permutation
of distinct ranks and every binary input. In the unbounded integer
encoding each comparator imposes $h=\max(x,y)$ and $l=x+y-h$.
First exclude a one-step increase greater than eight. For depth nine,
partition by the sixteen sign choices $4g_i\ge B$ or $4g_i<B$ at the
source. In seeking $V_9\ge V_0>200$, the first inequality permits adding
$V_0-8(9-j)\le V_j\le V_0+8j$ and $V_j\le V_{j-1}+8$ for $j\ge1$.
There is no upper bound on $B$ in any query. All seventeen negations
are UNSAT in Z3 and Yices. Two false strengthenings yield SAT, and their
paths are checked by ordinary digit subtraction.

\section*{Data availability}
The ancillary material is archived at Zenodo as version~1.5 of \href{https://doi.org/10.5281/zenodo.24325109}{doi:10.5281/zenodo.24325109}, and maintained at \url{https://github.com/galyavin/kaprekar-dynamics}.

\section*{Tools}
The research, computations and drafting were carried out with substantial
assistance from AI systems (OpenAI Codex and Anthropic Claude), under the
author's direction and close supervision. The main software was Python,
Z3, cvc5, Yices and a C++ compiler. The author takes responsibility for
the contents.

\begin{thebibliography}{HPSW}
\raggedright
\bibitem[COST]{COST} E. Chen, K. Ono, R. E. Schwartz and D. Thakur,
\emph{Four-digit Kaprekar dynamics in odd bases}, J. Integer Seq. \textbf{29}
(2026), Article 26.4.7; \href{https://arxiv.org/abs/2606.20439}{arXiv:2606.20439}.
\bibitem[G1]{G1} R. Galyavin, \emph{Cycles of the four-digit Kaprekar map in
even bases}, 2026, in \emph{Kaprekar dynamics: papers, verification code
and data}, v1.4, \href{https://doi.org/10.5281/zenodo.24343078}{doi:10.5281/zenodo.24343078}.
\bibitem[G2]{G2} R. Galyavin, \emph{Arrival times for the four-digit Kaprekar
map in even bases}, 2026, ibid., v1.1,
\href{https://doi.org/10.5281/zenodo.24329321}{doi:10.5281/zenodo.24329321}.
\bibitem[G3]{G3} R. Galyavin, \emph{A common arrival-time law for five-digit
Kaprekar dynamics}, 2026, ibid., v1.2,
\href{https://doi.org/10.5281/zenodo.24330159}{doi:10.5281/zenodo.24330159}.
\bibitem[G4]{G4} R. Galyavin, \emph{Majorization in odd-digit Kaprekar
dynamics and the nine-digit arrival law}, 2026, ibid., v1.3,
\href{https://doi.org/10.5281/zenodo.24342270}{doi:10.5281/zenodo.24342270}.
\bibitem[HPSW]{Het} E. Hetmaniok, M. Pleszczy\'nski, I. Sobstyl and
R. Witu{\l}a, \emph{Kaprekar's transformations. Part II---numerical results
and intriguing corollaries}, Position Papers of the Federated Conference on
Computer Science and Information Systems, ACSIS \textbf{6} (2015), 97--104,
\href{https://doi.org/10.15439/2015F15}{doi:10.15439/2015F15}.
\bibitem[Kah]{Kah} M. Kahramanlar, \emph{An ordered-gap simplex reduction
for Kaprekar dynamics in arbitrary digit length and base}, preprint (21 July
2026), \href{https://doi.org/10.5281/zenodo.21472305}{doi:10.5281/zenodo.21472305}.
\bibitem[KD22]{KD22} A. Kay and K. Downes-Ward, \emph{Fixed points and cycles
of the Kaprekar transformation: 1. Odd bases}, J. Integer Seq. \textbf{25}
(2022), Article 22.6.7.
\bibitem[KD24]{KD24} A. Kay and K. Downes-Ward, \emph{Fixed points and cycles
of the Kaprekar transformation: 2. Even bases}, preprint (2024),
\href{https://arxiv.org/abs/2408.12257}{arXiv:2408.12257}.
\bibitem[N]{Nu} F. Nuez, \emph{Symmetries in generalized Kaprekar's
routine}, Symmetry \textbf{14} (2022), no.~1, article 37,
\href{https://doi.org/10.3390/sym14010037}{doi:10.3390/sym14010037}.
\bibitem[T]{Th} D. S. Thakur, \emph{Kaprekar phenomena}, in: Topics in
Number Theory, RMS Lecture Notes Series \textbf{26}, Ramanujan Mathematical
Society, 2020, 61--70.
\bibitem[W]{Wa} B. L. Walden, \emph{Searching for Kaprekar's constants:
algorithms and results}, International Journal of Mathematics and
Mathematical Sciences \textbf{2005} (2005), no.~18, 2999--3004,
\href{https://doi.org/10.1155/IJMMS.2005.2999}{doi:10.1155/IJMMS.2005.2999}.
\end{thebibliography}
\end{document}
